\definecolor{blue}{RGB}{41,5,195}
\makeatletter
\hypersetup{
     	%pagebackref=true,
		pdftitle={\@title}, 
		pdfauthor={\@author},
    	pdfsubject={\imprimirpreambulo},
	    pdfcreator={LaTeX with abnTeX2},
		pdfkeywords={abnt}{latex}{abntex}{abntex2}{trabalho acadêmico}, 
		colorlinks=true,       		% false: boxed links; true: colored links
    	linkcolor=blue,          	% color of internal links
    	citecolor=blue,        		% color of links to bibliography
    	filecolor=magenta,      		% color of file links
		urlcolor=blue,
		bookmarksdepth=4
}
\makeatother
\setlength{\parindent}{1.3cm}
\setlength{\parskip}{0.2cm}  
\makeindex


% Prefixo do cabeçalho: "Capítulo" no corpo, "APÊNDICE" nos apêndices
\makeatletter
\providecommand{\headerchapname}{\chaptername}
\def\headerchapname{\chaptername}
\AtBeginEnvironment{apendicesenv}{\def\headerchapname{\MakeUppercase{\apendicename}}}
\makeatother

% Cabeçalho: "Capítulo N. Título" / "APÊNDICE A. Título", igual em todas as páginas
\makepsmarks{abntheadings}{%
	\def\chaptermark##1{\markboth{\headerchapname\ \thechapter.\ ##1}{}}%
	\def\sectionmark##1{}%
	\def\subsectionmark##1{}%
}
\makeevenhead{abntheadings}
{\ABNTEXfontereduzida\textit\leftmark}{}{\ABNTEXfontereduzida\thepage}
\makeoddhead{abntheadings}
{\ABNTEXfontereduzida\textit\leftmark}{}{\ABNTEXfontereduzida\thepage}